\documentclass[12pt,a4paper]{article}

\usepackage[T1]{fontenc}
\usepackage[utf8]{inputenc}
\usepackage[french]{babel}

\usepackage[top=2cm, bottom=2cm, left=2cm, right=2cm]{geometry}
\usepackage{amsmath}
\usepackage{amssymb}
\usepackage{array}
\usepackage{tabularx}
\usepackage{booktabs}
\usepackage{fancyhdr}
\usepackage{lastpage}
\usepackage{mdframed}
\usepackage{xcolor}
\usepackage{enumitem}
\usepackage[version=4]{mhchem}          % formules chimiques : \ce{H2O}, \ce{CO2}
\usepackage{hyperref}
\usepackage{caption}
\usepackage{float}
\usepackage{tikz}            % schémas simples si besoin
\usepackage{pgfplots}
\pgfplotsset{compat=1.18}

% ============================================================
%  EN-TÊTE / PIED DE PAGE
% ============================================================
\pagestyle{fancy}
\fancyhf{}
\renewcommand{\headrulewidth}{0pt}

% Pied de page : référence sujet à gauche, numéro de page à droite
\fancyfoot[L]{\small BB26PC1}
\fancyfoot[R]{\small Page \thepage\ sur \pageref{LastPage}}

% ============================================================
%  COMMANDES UTILITAIRES
% ============================================================
% Titre de question souligné : \question{Titre}
\newcommand{\question}[1]{\par\medskip\noindent\underline{\textbf{#1}}\par\smallskip}

% Sous-question numérotée : \sousq{2a}{Texte de la question}
\newcommand{\sousq}[2]{\par\noindent\textbf{#1-} #2\par\smallskip}

% Encadré "Document" : \begin{document_box}{Titre} ... \end{document_box}
\newmdenv[
  linewidth=0.8pt,
  innerleftmargin=6pt,
  innerrightmargin=6pt,
  innertopmargin=4pt,
  innerbottommargin=4pt,
]{document_box}

% Points d'une question entre parenthèses
\newcommand{\pts}[1]{\textbf{(#1 points)}}

% ============================================================
%  DÉBUT DU DOCUMENT
% ============================================================
\begin{document}

% ============================================================
%  PAGE DE GARDE (première page)
% ============================================================
\thispagestyle{fancy}

\begin{center}
  {\large\bfseries BREVET BLANC}\\[30pt]
  {\large SESSION 2026}\\[30pt]

  \fbox{\parbox{0.92\linewidth}{%
    \begin{center}
      {\Large\bfseries PHYSIQUE-CHIMIE}\\[20pt]
      \textit{Durée de l'épreuve : 30 minutes} – \textbf{20 points}
    \end{center}
  }}

\vspace{30pt}
\noindent
Dès que le sujet vous est remis, assurez-vous qu'il est complet.\\[20pt]
Ce sujet comporte 3 pages numérotées de la page 1/3 à la page 3/3.\\[20pt]
L’utilisation du dictionnaire n’est pas autorisée.\\
L’usage de la calculatrice sans mémoire « type collège » est autorisé.\\[30pt]
Les essais et les démarches engagés, même non aboutis, seront pris en compte.
\end{center}

\vspace{1cm}

% ============================================================
%  PARTIE CHIMIE
% ============================================================
\newpage

\section*{Partie 1 : Activités humaines et conséquences }

%  --- Texte d'introduction ---
\noindent
Depuis plus de 260 ans, avec le début de la révolution industrielle, les activités humaines n’ont cessé d’augmenter entraînant la production et le rejet massif de gaz à effet de serre dans l’atmosphère.\\
La production d’électricité est la première cause d’émission de gaz à effet de serre. Parmi ces derniers, c’est le dioxyde de carbone (\ce{CO2}) qui est le plus produit, notamment lors de la combustion des ressources fossiles. L’augmentation des concentrations atmosphériques en \ce{CO2} et autres gaz à effet de serre est responsable d’importants changements climatiques. Outre le réchauffement climatique, l’augmentation du taux de dioxyde de carbone dans l’atmosphère provoque également une acidification des océans et modifie la croissance des végétaux.

\vspace{0.4cm}

% --- Question 1 PC ---
\question{Question 1 \pts{4}}

\vspace{0.4cm}

\sousq{1a}{\textbf{Donner} la cause principale de l’augmentation du taux de dioxyde de carbone (\ce{CO2}) dans l’atmosphère, d’après le texte d’introduction ci-dessus.}

\vspace{0.4cm}

\sousq{1b}{\textbf{Donner} trois conséquences de l’augmentation du taux de dioxyde de carbone dans l’atmosphère, toujours d’après le texte d’introduction.}

\vspace{0.4cm}

% --- Document 1 ---
\fbox{\parbox{0.95\linewidth}{%
\question{Document 1 : Acidification des océans}
\noindent
Le dioxyde de carbone (\ce{CO2}) réagit avec l’eau (\ce{H2O}) des océans suivant la transformation chimique suivante :
\[ \ce{CO2 + H2O -> HCO3- + H+} \]\\[1pt]
L’augmentation du taux de \ce{CO2} sur Terre entraîne donc à son tour une augmentation de la concentration en ions hydrogène \ce{H+}. Cette augmentation est responsable d’une modification du pH de l’eau de mer. L’eau de mer devient plus acide : c’est ce qu’on appelle l’acidification des océans.\\[2pt]
Avant la survenue de ce phénomène, la valeur du pH moyen des océans était de 8,2. 
}}

\vspace{0.4cm}

% --- Question 2 ---
\question{Question 2 \pts{3}}

\vspace{0.4cm}

\noindent
\textbf{Recopier et compléter} les phrases suivantes avec la bonne proposition.\\[1pt]

\sousq{2a}{Si le taux de CO2 continue d’augmenter, à l’avenir, le pH moyen des océans sera :}
\begin{center}
  $\triangleright$ inférieur à 8,2
  \hspace{2cm}
  $\triangleright$ égal à 8,2
  \hspace{2cm}
  $\triangleright$ supérieur à 8,2
\end{center}

\textbf{Justifier} la réponse.

\vspace{0.4cm}

\sousq{2b}{Le matériel qui permet de mesurer l’acidité d’une solution est le :}
\begin{center}
  $\triangleright$ thermomètre
  \hspace{2cm}
  $\triangleright$ dynamomètre
  \hspace{2cm}
  $\triangleright$ papier pH
  \hspace{2cm}
  $\triangleright$ voltmètre
\end{center}

\vspace{0.4cm}

% --- Question 3 ---
\question{Question 3 \pts{2}}

\vspace{0.4cm}

\noindent
\textbf{Donner} le nom et le nombre d’atomes présents dans la molécule d'eau.

% ============================================================
%  PARTIE MOUVEMENT
% ============================================================

\newpage

\section*{Partie 2 : Tirs au but}

%  --- Texte d'introduction ---
\noindent
Désiré et Matvei s'entraînent lors d'une séance de tirs au but au football. Depuis plusieurs années, la technologie permet de suivre la trajectoire des ballons à l'aide de capteurs. Ces outils aident à analyser les matchs et à optimiser les entraînements, notamment en calculant la vitesse du ballon.

\vspace{0.4cm}

\question{Question 4 \pts{2}}

\vspace{0.4cm}

\noindent
\textbf{Donner} la formule permettant de calculer une vitesse moyenne. Préciser le nom et le symbole de chaque grandeur qui y apparaît.\\[0.4cm]
Matvei envoie le ballon à Désiré, qui se trouve au milieu du terrain. Le ballon parcourt 60,8 m en 3,2 s.

\vspace{0.4cm}

\question{Question 5 \pts{2}}

\vspace{0.4cm}

\noindent
\textbf{Calculer} la vitesse moyenne du ballon. On détaillera le calcul et on précisera l'unité du résultat.

\vspace{0.4cm}

\question{Document 2 : Vitesse moyenne de la balle au cours du temps}

\begin{tikzpicture}
\begin{axis}[
    xlabel={\textbf{Temps} (s)},
    ylabel={\textbf{Vitesse moyenne} (m/s)},
    xmin=0, xmax=3.5,
    ymin=0, ymax=30,
    xtick={0, 0.4, 0.8, 1.2, 1.6, 2.0, 2.4, 2.8, 3.2},
    ytick={0, 5, 10, 15, 20, 25, 30, 35, 40, 45},
    grid=both,
    grid style={line width=0.2pt, draw=gray!30},
    major grid style={line width=0.4pt, draw=gray!60},
    width=10cm, height=7cm,
    mark size=2pt,
]
\addplot[
    color=black,
    mark=*,
    thick,
] coordinates {
    (0, 27)
    (0.4, 26.2)
    (0.8, 24.3)
    (1.2, 23)
    (1.6, 22)
    (2.0, 20)
    (2.4, 16)
    (2.8, 12)
    (3.2, 8)
};
\end{axis}
\end{tikzpicture}

\question{Question 6 \pts{2}}

\vspace{0.4cm}

\noindent
En vous appuyant sur le Document 2, \textbf{indiquer} si le mouvement du ballon est accéléré, uniforme ou ralenti.\\[0.4cm]
Désiré se place sur le point de penalty, à 11 m du but. Il vise la lucarne : la distance que doit parcourir le ballon est alors de 11,50 m. Lors de sa frappe, la vitesse moyenne du ballon est de 30 m/s. Le temps de réaction moyen d'un être humain est de 0,4 s.

\vspace{0.4cm}

\question{Question 7 \pts{5}}

\vspace{0.4cm}

\noindent
\textbf{Expliquer} pourquoi Matvei doit anticiper la frappe de Désiré s'il veut avoir une chance d'arrêter le ballon. Un raisonnement s’appuyant sur des calculs est attendu.


\end{document}